\exercisesection{1}

\begin{enumerate}
\def\labelenumi{\arabic{enumi}.}
\tightlist
\item
  \begin{enumerate}
  \def\labelenumii{(\alph{enumii})}
  \tightlist
  \item
    设 A 是绝对平坦环（第 I 章，§2，习题 17）。证明 Ass(A) 是 Spec(A) 的孤立点组成的集（注意 Ass(A) 是 A 的幂等元的零化子这样的素理想组成的集）。
  \end{enumerate}
\end{enumerate}

\begin{enumerate}
\def\labelenumi{(\alph{enumi})}
\setcounter{enumi}{1}
\item
  由 (a) 推出一个环 A 的例子，使 \(A \neq 0\) 且 \(\operatorname{Ass}(A) = \varnothing\) （参见第 II 章，§4，习题 17），并且第 1 节命题 2 的推论 2 的结论对它不成立。
\item
  由 (b) 推出一个环 A 及 A 的乘性子集 S 的例子，使得 \(p \mapsto S^{-1}p\) 这一从 \(\text{Ass}(A) \cap \Phi\) 到 \(\text{Ass}(S^{-1}A)\) 的映射（第 2 节，命题 5）不是满射。
\end{enumerate}

* 2. 设 K 是交换域，A 是以 Q 为阶群的赋值环，由``形式幂级数'' \(\sum_{r} c_r T^r\) 组成，其中 \(c_r \in K\) ，\(r \in \mathbf{Q}_+\) ，且诸 \(r\) 中使 \(c_r \neq 0\) 者所成的集是良序的。设 \(\alpha\) 是无理数 \(>0\) ，\(a\) 是 A 中由赋值 \(>\alpha\) 的元素组成的理想，C 是环 A/a。那么 \(\text{Spec(C)}\) 只剩一点 \(p\) ；对 C 中所有 \(x \neq 0\) ，有 \(px \neq 0\) ，但对所有 \(\lambda \in p\) ，存在 C 中的 \(y \neq 0\) 使 \(\lambda y = 0\) 。特别地，\(\text{Ass(C)} = \varnothing\) ，尽管 \(\text{Supp(C)} = \text{Spec(C)} = \{\mathfrak{p}\}\) 。 *

\begin{enumerate}
\def\labelenumi{\arabic{enumi}.}
\setcounter{enumi}{2}
\item
  设 M 是 A-模，N 是 M 的子模。证明每个不包含 Ann(N) 的素理想 \(\mathfrak{p} \in \operatorname{Ass}(M/N)\) 都伴随于 M。给出一个属于 Ass(M/N) 而不属于 Ass(M) 的理想 p 的例子（取 A 为整环，M = A）。
\item
  给出一个环 A 的例子，使 M = A 不满足第 4 节定理 1 的结论（参见习题 1 (b)）。
\item
  设 A = K{[}X, Y{]} 是交换域 K 上两个未定元的多项式代数，a 是 A 的极大理想 AX + AY。证明 \(\text{Supp}(a) = \text{Spec}(A)\) 是无限集，且 \(\text{Ass}(a) = \{0\}\) 。证明 a 没有所有因子模都与 A 同构的合成列。（这样的列的存在将意味着 a 同构于某个模 \(A^{n}\) ；于是必有 n = 1，这是荒谬的。）
\end{enumerate}

¶ 6. 设 A 是环，E 是 A-模。

\begin{enumerate}
\def\labelenumi{(\alph{enumi})}
\item
  证明：要使 A 的素理想 p 属于 Supp(E)，必须且只需存在 E 的子模 F 使 \(p \in \text{Ass}(E/F)\) （要证条件必要，考虑 E 中形如 px 的子模 F，其中 \(x \in E\) ；要证条件充分，用第 3 节命题 7）。
\item
  假设 A 是 Noether 环且 E 有限生成。证明对每个素理想 \(\mathfrak{p}\in\operatorname{Supp}(\mathrm{E})\) ，存在 E 的合成列 \((\mathrm{E}_{i})_{0\leqslant i\leqslant n}\) ，使得对 \(0\leqslant i\leqslant n-1\) ，\(E_{i}/E_{i+1}\) 同构于 \(A/p_{i}\) ，其中 \(p_{i}\) 是素理想，且诸 \(p_{i}\) 中有一个等于 p。
\end{enumerate}

\begin{enumerate}
\def\labelenumi{\arabic{enumi}.}
\setcounter{enumi}{6}
\tightlist
\item
  设 A 是 Noether 环，M 是有限生成 A-模，\(\mathfrak{a} = \text{Ann}(M)\) 。
\end{enumerate}

\begin{enumerate}
\def\labelenumi{(\alph{enumi})}
\item
  证明：若 \(\mathfrak{a}\) 是素的，则 \(\mathfrak{a}\) 是 \(\operatorname{Ass}(\mathbf{M})\) 的最小元。（注意 \(\mathfrak{a} = \bigcap_{\mathfrak{p} \in \operatorname{Ass}(\mathbf{M})} \mathfrak{p}\) 。）
\item
  证明伴随于 A/a 的每个素理想都伴随于 M。（注意若 p 是 A 的素理想，且是元素 \(\alpha \in A\) 在 A/a 中的类的零化子，则 \(\mathfrak{p} = \operatorname{Ann}(\alpha\mathbf{M})\) ，然后用 (a)。）
\item
  设 \(p, q\) 是 \(A\) 的两个不同的素理想，满足 \(p \subset q\) 。证明：若 \(M = (A/p) \oplus (A/q)\) ，则 \(\operatorname{Ass}(A/a) \neq \operatorname{Ass}(M)\) 。
\end{enumerate}

\begin{enumerate}
\def\labelenumi{\arabic{enumi}.}
\setcounter{enumi}{7}
\item
  设 A 是 Noether 环，a 是 A 的理想，M 是有限生成 A-模，P 是 M 中由满足 ax = 0 的 \(x \in M\) 组成的子模。证明 \(\operatorname{Ass}(M/P) \subset \operatorname{Ass}(M)\) （注意对每个 \(x \in M\) ，\((Ax + P)/P\) 的零化子也是 ax 的零化子，并用习题 7 (a)）。
\item
  设 A 是 Noether 环，a 是 A 的理想。
\end{enumerate}

\begin{enumerate}
\def\labelenumi{(\alph{enumi})}
\item
  要使 A 的理想 b 满足 a: b ≠ a，必须且只需 b 含于某个素理想 p ∈ Ass(A/a)。
\item
  设 A 是域 K 上 3 个未定元的多项式代数 K{[}X, Y, Z{]}。设 n = AX，m = AX + AY + AZ，a = n ∩ m²。证明存在包含 a 的素理想 p 使 a: p ≠ a，但它不伴随于 A。
\end{enumerate}

\begin{enumerate}
\def\labelenumi{\arabic{enumi}.}
\setcounter{enumi}{9}
\tightlist
\item
  \begin{enumerate}
  \def\labelenumii{(\alph{enumii})}
  \tightlist
  \item
    给出一个 Z-模 M 的例子，使 \(\operatorname{Ass}(\operatorname{Hom}_{\mathbf{Z}}(\mathbf{M}, \mathbf{M}))\) 不含于 \(\operatorname{Supp}(\mathbf{M})\) （参见第 II 章，§ 4，习题 24 (c)）。
  \end{enumerate}
\end{enumerate}

\begin{enumerate}
\def\labelenumi{(\alph{enumi})}
\setcounter{enumi}{1}
\tightlist
\item
  给出 Z-模 E、F 的例子，使得
\end{enumerate}

\[
\operatorname{Ass} \left(\operatorname{Hom} _ {\mathbf {Z}} (\mathrm{E}, \mathrm{F})\right) = \varnothing ,
\]

但 \(\operatorname{Ass}(\mathbf{F}) \cap \operatorname{Supp}(\mathbf{E}) \neq \varnothing\) （取 \(\mathbf{F} = \mathbf{Z}\) ）。

\begin{enumerate}
\def\labelenumi{\arabic{enumi}.}
\setcounter{enumi}{10}
\tightlist
\item
  \begin{enumerate}
  \def\labelenumii{(\alph{enumii})}
  \tightlist
  \item
    设 \(\mathbf{A}\) 是 Noether 环，\(\mathbf{E}\) 是 \(\mathbf{A}\) -模。证明从 E 到乘积 \(\prod_{p\in\text{Ass}(E)}E_p\) 的典范同态是单射（若 N 是这个同态的核，证明 \(\text{Ass}(N)=\varnothing\) ）。
  \end{enumerate}
\end{enumerate}

\begin{enumerate}
\def\labelenumi{(\alph{enumi})}
\setcounter{enumi}{1}
\tightlist
\item
  取 A 为域 K 上的多项式环 K{[}X,Y,Z{]}；设 \(p_{1}=AX+AY\) ，\(p_{2}=AX+AZ\) ，它们是素理想，又设 a 为理想 \(p_{1}p_{2}\) 。集合 \(\operatorname{Ass}(A/a)\) 由 \(p_{1}\) 、\(p_{2}\) 及极大理想 \(m=p_{1}+p_{2}\) 组成；证明从 E=A/a 到 \(E_{p_{1}}\times E_{p_{2}}\) 的典范同态不是单射。
\end{enumerate}

\begin{enumerate}
\def\labelenumi{\arabic{enumi}.}
\setcounter{enumi}{11}
\tightlist
\item
  设 A 是 Noether 环，P 是投射 A-模。证明：若对每个 \(p \in \text{Ass}(A)\) ，\(P_{p}\) 都是有限生成 \(A_{p}\) -模，则 P 是有限生成 A-模（把 A 嵌入乘积 \(\prod_{p \in \text{Ass}(A)} A_{p}\) （习题 11），并用《代数》第 II 章，§ 5，第 5 节，命题 9）。
\end{enumerate}

¶ 13. 设 A 是 Noether 环，E 是有限生成 A-模，p、q 是 A 的两个素理想，满足 \(p \subset q\) ，a 是 q 的一个元素。假设 \(p \in \text{Ass}(E)\) 且位似映射 \(a_{E}\) 是单射。证明存在素理想 \(n \in \text{Ass}(E/aE)\) 使 \(p + Aa \subset n \subset q\) 。（把 A 换成 \(A_{q}\) ，可设 A 是以 q 为极大理想的局部环。设 \(F \neq 0\) 是 E 中由满足 px = 0 的 x 组成的子模。证明关系 \(F \subset aE\) 将导出 F = aF，并用 Nakayama 引理得出矛盾。然后用第 4 节命题 8。）

\begin{enumerate}
\def\labelenumi{\arabic{enumi}.}
\setcounter{enumi}{13}
\item
  设 A 是整环，\(a \neq 0\) 是 A 的元素，p 是 \(\operatorname{Ass}(A/aA)\) 的元素。证明对 p 的每个元素 \(b \neq 0\) ，有 \(\mathfrak{p} \in \operatorname{Ass}(A/bA)\) （若 \(c \in A\) 使关系 \(xc \in aA\) 等价于 \(x \in p\) ，证明存在 \(d \in A\) 使 \(xd \in bA\) 等价于 \(xc \in aA\) ）。
\item
  设 A 是 Noether 环，E 是有限生成 A-模，F 是 E 的子模，m 是 A 的理想。要使 F 在 E 中对于 m-进拓扑是闭的，必须且只需 \(p + m \neq A\) 对每个理想 \(p \in \text{Ass}(E/F)\) 成立。（化归到 F = 0 的情形，用第 III 章，§3，第 2 节，命题 5 的推论，并把第 1 节命题 2 的推论 2 应用于元素 \(1 + m\) ，其中 \(m \in m\) 。）
\item
  设 A 是 Noether 环。要使 A 同构于有限多个整环的乘积，必须且只需对 A 的每个素理想 p，局部环 \(A_{p}\) 都是整环。（要证条件充分，首先注意它蕴涵环 A 是约化的；由此推出 \(\{0\} = \bigcap_{i} p_{i}\) ，其中 \(p_{i} (1 \leqslant i \leqslant n)\) 是 A 的极小素理想。证明对 \(i \neq j\) ，必有 \(p_{i} + p_{j} = A\) ；为此注意，若存在包含 \(p_{i} + p_{j}\) 的极大理想 m，则环 \(A_{m}\) 就不是整环，这里要用第 1 节命题 2 的推论 3 及第 2 节命题 5 的推论。）
\end{enumerate}

¶ 17. 设 A 是环，M 是 A-模。若存在 \(x \in \mathbf{M}\) 使 p 是包含 Ann(x) 的素理想组成的集的极小元，则称 A 的素理想 p 弱伴随于 M；我们用 \(\operatorname{Ass}_{f}(\mathbf{M})\) 表示弱伴随于 M 的理想组成的集。那么 \(\operatorname{Ass}(\mathbf{M}) \subset \operatorname{Ass}_{f}(\mathbf{M})\) 。

\begin{enumerate}
\def\labelenumi{(\alph{enumi})}
\item
  证明关系 \(\mathbf{M} \neq 0\) 等价于 \(\operatorname{Ass}_f(\mathbf{M}) \neq \varnothing\) 。
\item
  要使 \(a \in \mathbf{A}\) 满足 \(a_{\mathbf{M}}\) 是单射，必须且只需 \(a\) 不属于 \(\operatorname{Ass}_f(\mathbf{M})\) 的任何元素（注意，若 \(a\) 属于某个理想 \(\operatorname{Ann}(x)\) 的根，其中 \(x \neq 0\) 在 \(\mathbf{M}\) 中，则存在 \(y \neq 0\) 于 \(\mathbf{M}\) 中，使 \(a \in \operatorname{Ann}(y)\) 。要证条件必要，考虑环 \(\operatorname{A}/\operatorname{Ann}(x)\) ，把它化为证明：在环 \(\mathbf{A}\) 中，属于素理想 \(p\) （它在 \(\mathbf{A}\) 的素理想组成的集中极小）的每个元素必是零因子（第 II 章，§ 2，第 6 节，命题 12）。）
\end{enumerate}

要使 \(a \in A\) 满足 \(a_{M}\) 是殆幂零的，必须且只需 a 属于 \(\operatorname{Ass}_{f}(M)\) 的所有元素。

\begin{enumerate}
\def\labelenumi{(\alph{enumi})}
\setcounter{enumi}{2}
\tightlist
\item
  若 \(\mathbf{N}\) 是 \(\mathbf{M}\) 的子模，则
\end{enumerate}

\[
\operatorname{Ass} _ {f} (\mathbf {N}) \subset \operatorname{Ass} _ {f} (\mathbf {M}) \subset \operatorname{Ass} _ {f} (\mathbf {N}) \cup \operatorname{Ass} _ {f} (\mathbf {M} / \mathbf {N}).
\]

（注意，若 \(\mathfrak{p}\) 是素理想，\(a \notin \mathfrak{p}\) ，\(x \in \mathbf{M}\) 满足 \(ax \in \mathbf{N}\) 且 \(\operatorname{Ann}(x) \subset \mathfrak{p}\) ，则 \(\operatorname{Ann}(ax) \subset \mathfrak{p}\) 。）

\begin{enumerate}
\def\labelenumi{(\alph{enumi})}
\setcounter{enumi}{3}
\item
  设 S 是 A 的乘性子集，\(\Phi\) 是 A 中与 S 不相交的素理想组成的集。证明 \(p \mapsto S^{-1}p\) 是从 \(\operatorname{Ass}_{f}(M) \cap \Phi\) 到 \(\operatorname{Ass}_{f}(S^{-1}M)\) 上的双射。（注意，在典范映射 \(A \to S^{-1}A\) 下，元素 x/s（其中 \(x \in M\) ，\(s \in S\) ）的零化子的原像是 \(\operatorname{Ann}(x)\) 关于 S 的饱和。）
\item
  设 S 是 A 的乘性子集；证明：若 N 是典范同态 \(M \to S^{-1}M\) 的核，则 \(\operatorname{Ass}_{f}(N)\) 是由 \(p \in \operatorname{Ass}_{f}(M)\) 中与 S 相交者组成的集，而 \(\operatorname{Ass}_{f}(M/N)\) 是由 \(p \in \operatorname{Ass}_{f}(M)\) 中与 S 不相交者组成的集。（为证后一点，考虑在包含 \(\operatorname{Ann}(\bar{y})\) 的理想中极小的素理想 q，其中 \(\bar{y} \in M/N\) ；注意，若 \(y \in \bar{y}\) 且 \(t \in q\) ，则存在 \(c \notin q\) 及整数 n \textgreater{} 0 使 \(ct^{n}y \in N\) （第 II 章，§2，第 6 节，命题 12）；先由此推出 \(q \cap S = \varnothing\) 。存在 \(s \in S\) 使 \(sct^{n}y = 0\) ；由此断定不可能有素理想 \(q' \neq q\) 使 \(\operatorname{Ann}(y) \subset q' \subset q\) 。）
\item
  要使 \(\mathbf{A}\) 的素理想属于 \(\operatorname{Supp}(\mathbf{M})\) ，必须且只需它包含 \(\operatorname{Ass}_f(\mathbf{M})\) 的一个元素。
\item
  证明：若 A 是 Noether 环，则 \(\operatorname{Ass}_f(M) = \operatorname{Ass}(M)\) 。
\item
  若 \(\mathbf{A}\) 是绝对平坦环，则 \(\operatorname{Ass}_f(\mathbf{A}) = \operatorname{Spec}(\mathbf{A})\) 。
\item
  对每个 A-模 E，从 E 到乘积
\end{enumerate}

\(\prod_{p\in\mathrm{Ass}_{f}(\mathbf{E})}\mathrm{E}_{p}\) 的典范同态是单射；换言之，\(\{0\}\) 在 E 中关于诸理想 \(p\in\mathrm{Ass}_{f}(\mathbf{E})\) 的饱和的交化为 0。类似地推广习题 12。

\begin{enumerate}
\def\labelenumi{(\alph{enumi})}
\setcounter{enumi}{9}
\tightlist
\item
  设 M 是有限生成 A-模。证明包含 \(\mathfrak{a} = \text{Ann}(M)\) 的素理想组成的集的极小元都属于 \(\text{Ass}_{f}(M)\) （若 \((x_{i})_{1 \leq i \leq n}\) 是 M 的生成系，证明这样的理想包含某个 \(\text{Ann}(x_{i})\) ）。由此推出 \(\text{Ass}_{f}(A/a)\) 含于 \(\text{Ass}_{f}(M)\) （若素理想 p 包含元素 \(\alpha \in A\) 模 a 的类的零化子，注意 p 包含 M 的子模 \(\alpha M\) 的零化子）；证明 \(\text{Ass}_{f}(A/a)\) 、\(\text{Ass}(M)\) 与包含 a 的素理想组成的集有相同的极小元。
\end{enumerate}

\begin{enumerate}
\def\labelenumi{\arabic{enumi}.}
\setcounter{enumi}{17}
\tightlist
\item
  设 A 是环，M 是 A-模。要使元素 \(c \in A\) 满足：对每个 \(a \in A\) （其位似映射 \(a_{M}\) 为单射），\(b_{M}\) 对所有形如 \(a + \lambda c\) （其中 \(\lambda \in A\) ）的 b 都是单射位似映射，只需 c 属于 \(\operatorname{Ass}_{f}(M)\) 的极大元之交（用习题 17 (b)）。这个条件必要吗？
\end{enumerate}

* 19. 设 A 是高为 2 的赋值的环，\(\Gamma\) 是它的阶群，\(\Gamma_1\) 是 \(\Gamma\) 的异于 \(\{0\}\) 与 \(\Gamma\) 的唯一孤立子群，\(\gamma > 0\) 是 \(\Gamma\) 中不属于 \(\Gamma_1\) 的元素。设 \(a\) 是由 \(x \in A\) 中使 \(v(x) \geqslant \gamma\) 的元素组成的理想，\(b\) 是由 \(x \in A\) 中使 \(v(x) > \gamma\) 的元素组成的理想，\(E = A/a\) ，\(F = A/b\) 。证明 \(\operatorname{Ass}_f(\operatorname{Hom}_A(E, F))\) 异于 \(\operatorname{Ass}_f(F) \cap \operatorname{Supp}(E)\) 。 *

